\subsection{Cohen rings}

DEFINITION 2. --- Let A be a separated and complete local ring, whose residue field is of characteristic p. One calls a Cohen subring of A a subring C of A which is a p-ring such that \(\mathbf{A} = \mathfrak{m}_{\mathbf{A}} + \mathbf{C}\) (i.e. \(\mathbf{A}/\mathfrak{m}_{\mathbf{A}} = \mathbf{C}/(\mathfrak{m}_{\mathbf{A}} \cap \mathbf{C})\) ).

If C is a Cohen subring of A, the ideal \(m_{A} \cap C\) of C is maximal, hence equal to pC. The canonical map from \(\kappa_{C} = C/pC\) onto \(\kappa_{A} = A/m_{A}\) is therefore an isomorphism of fields.

Example. --- Let C be a p-ring. The ring of formal power series \(A = C[[T_{1}, \ldots, T_{n}]]\) is a Noetherian local ring, separated and complete, whose maximal ideal is generated by the sequence (p, T₁, \ldots, Tₙ). It is immediate that C is a Cohen subring of A. This applies in particular when C is equal to Zₚ, to Z/pⁿZ, or to a field of characteristic p.

THEOREM 1. — Let A be a local ring, separated and complete, whose residue field k has characteristic p. Let π be the canonical map from A onto k, and let S be a subset of A such that π induces a bijection from S onto a p-basis of k (A, V, p. 95).

\begin{enumerate}
\def\labelenumi{\alph{enumi})}
\item
  There exists a Cohen subring C of A containing S, and it is unique.
\item
  The subring C of A is closed, and the \(p\)C-adic topology on C is induced by the \(m_{A}\)-adic topology of A.
\item
  Every closed subring A' of A, containing S, and such that \(A = A' + m_{A}\) contains C.
\end{enumerate}

\begin{enumerate}
\def\labelenumi{\Alph{enumi})}
\tightlist
\item
  Special case: \(\mathfrak{m}_{\mathbf{A}}\) nilpotent
\end{enumerate}

Let \(n\) be a positive integer such that \(\mathfrak{m}_{\mathbf{A}}^{n+1}=\{0\}\) . If \(\Phi_{n}\) is the \(n\) -th Witt polynomial ( § 1, no. 1 ), the map \(u:[a_{0},..., a_{n}]\mapsto\Phi_{n}(a_{0},..., a_{n})\) is a ring homomorphism from \(\mathbf{W}_{n+1}(\mathbf{A})\) to \(\mathbf{A}\) ( § 1, no. 7 ). Let \(\mathbf{B}_{n}\) be the image of \(u\) and let \(\mathbf{C}_{n}\) be the subring of \(\mathbf{A}\) generated by \(\mathbf{B}_{n}\cup\mathbf{S}\) .

Lemma 1. --- Let A' be a subring of A containing S. For A' to contain \(\mathbf{C}_{n}\) , it is necessary and sufficient that \(A' + m_{A} = A\).

We have \(pA \subset m_A\) and \(B_n\) consists of the elements of the form \(a_0^{p^n} + pa_1^{p^{n-1}} + \cdots + p^n a_n\) with \(a_0, \ldots, a_n\) in A. Consequently, we have \(\pi(B_n) = k^{p^n}\) , whence \(\pi(C_n) = k^{p^n}[\pi(S)]\) . But since \(\pi(S)\) is a \(p\) -basis of \(k\) , one has \(k = k^{p^n}[\pi(S)]\) (A, V, p. 96), whence \(\pi(C_n) = k\) , that is, \(C_n + m_A = A\) .

Let A' be a subring of A containing S. If A' contains \(C_{n}\) , one has

\[
\mathrm{A} ^ {\prime} + \mathrm{m} _ {\mathrm{A}} \supset \mathrm{C} _ {n} + \mathrm{m} _ {\mathrm{A}} = \mathrm{A}, \quad \mathrm{whence} \quad \mathrm{A} ^ {\prime} + \mathrm{m} _ {\mathrm{A}} = \mathrm{A}.
\]

Conversely, suppose we have \(\mathbf{A}' + \mathfrak{m}_{\mathbf{A}} = \mathbf{A}\) . Let \(a_0, \ldots, a_n\) be elements of \(\mathbf{A}\) ; by hypothesis there exist elements \(a_0', \ldots, a_n'\) of \(\mathbf{A}'\) such that \(a_i \equiv a_i'\) mod. \(\mathfrak{m}_{\mathbf{A}}\) for \(0 \leqslant i \leqslant n\) . By prop. 1 of § 1, no. 1 and the hypothesis \(\mathfrak{m}_{\mathbf{A}}^{n+1} = \{0\}\) , one therefore has \(\Phi_n(a_0, \ldots, a_n) = \Phi_n(a_0', \ldots, a_n') \in \mathbf{A}'\) , whence \(\mathbf{B}_n \subset \mathbf{A}'\) . Since \(\mathbf{C}_n\) is the ring generated by \(\mathbf{B}_n \cup \mathbf{S}\) , one has \(\mathbf{C}_n \subset \mathbf{A}'\) .

In the set S of subrings A' of A containing S and such that \(A' + m_{A} = A\) there exists, by lemma 1, a smallest element C, and one has \(C_{n} = C\) for every integer \(n \geqslant 0\) such that \(m_{A}^{n+1} = \{0\}\) .

We have \(C + m_A = A\) by construction and \(p1_C\) is nilpotent. One obviously has \(pC \subset C \cap m_A\) and lemma 2 below therefore shows that \(pC\) is a maximal ideal of \(C\) and consequently that \(C\) is a Cohen subring of \(A\) .

Lemma 2. --- One has \(C \cap m_{A} \subset pC\).

Let us choose an integer \(m \geqslant 1\) such that \(\mathfrak{m}_{\mathbf{A}}^{m} = \{0\}\) , whence \(\mathbf{C} = \mathbf{C}_{m} = \mathbf{C}_{m-1}\) . Let \(\Lambda\) be the part of \(\mathbf{N}^{(\mathbf{S})}\) consisting of families with finite support of integers \((\alpha_{s})_{s \in \mathbf{S}}\) satisfying \(0 \leqslant \alpha_{s} < p^{m}\) for every \(s \in \mathbf{S}\) . Since \(\mathbf{B}_{m}\) contains \(s^{p^{m}} = \Phi_{m}(s, 0, ..., 0)\) for every \(s \in \mathbf{S}\) , the monomials \(Z_{\alpha} = \prod_{s \in \mathbf{S}} s^{\alpha_{s}}\) , where \(\alpha\) ranges over \(\Lambda\) , generate \(\mathbf{C}_{m}\) as a \(\mathbf{B}_{m}\) -module. Moreover, by the formula

\[
\Phi_ {m} (a _ {0}, \dots , a _ {m}) = a _ {0} ^ {p ^ {m}} + p \Phi_ {m - 1} (a _ {1}, \dots , a _ {m}),
\]

every element of \(B_{m}\) is of the form \(a^{p^{m}} + pb\) with \(a \in A\) and \(b \in B_{m-1}\) . Consequently every element of \(C = C_{m}\) is of the form

\[
x = \sum_ {\alpha \in \Lambda} c _ {\alpha} ^ {p ^ {m}} Z _ {\alpha} + p y\tag{1}
\]

with \(c_{\alpha} \in \mathbf{A}\) for every \(\alpha \in \Lambda\) , and \(y \in \mathbf{C}_{m-1} = \mathbf{C}\) . If \(x\) belongs to \(\mathbf{C} \cap \mathfrak{m}_{\mathbf{A}}\) , one has \(\pi(x) = 0\) , whence \(\sum_{\alpha \in \Lambda} \pi(c_{\alpha})^{p^m} \pi(Z_{\alpha}) = 0\) . Since \(\pi(\mathbf{S})\) is a \(p\) -basis of \(k\) , one has \(\pi(c_{\alpha}) = 0\) for every \(\alpha \in \Lambda\) according to A, V, p. 96. We then have \(c_{\alpha} \in \mathfrak{m}_{\mathbf{A}}\) , whence \(c_{\alpha}^m = 0\) and a fortiori \(c_{\alpha}^{p^m} = 0\) . By (1), we have \(x = py\) , whence lemma 2.

We have \(p^m\mathbf{C} = \mathfrak{m}_{\mathbf{A}}^m = \{0\}\) for \(m\) sufficiently large, and assertion \(b)\) is therefore trivial. Assertion \(c)\) follows from lemma 1. If \(\mathbf{C}'\) is a Cohen subring of A containing S, we have \(\mathbf{C}' \supset \mathbf{C}\) by lemma 1. But since the inclusion of C in \(\mathbf{C}'\) induces an isomorphism of \(\kappa_{\mathbf{C}}\) onto \(\kappa_{\mathbf{C}'}\) , one has \(\mathbf{C} = \mathbf{C}'\) (no. 1, prop. 2, b)), and this completes the proof of \(a\) ).

\noindent\textbf{B) General case.}\par

For every integer \(n \geqslant 0\) , denote by \(\mathbf{A}_n\) the local ring \(\mathbf{A}/\mathfrak{m}_{\mathbf{A}}^{n+1}\) , \(\mathfrak{m}_n = \mathfrak{m}_{\mathbf{A}}/\mathfrak{m}_{\mathbf{A}}^{n+1}\) its maximal ideal, and \(\pi_n\) the canonical homomorphism from A onto \(\mathbf{A}_n\) . By A), there exists a unique Cohen subring \(\mathbf{C}_n\) of \(\mathbf{A}_n\) containing \(\pi_n(\mathbf{S})\) . When \(0 \leqslant n \leqslant m\) , denote by \(\pi_{n,m}\) the canonical homomorphism from \(\mathbf{A}_m\) onto \(\mathbf{A}_n\) . By cor. 2 of prop. 1 of no. 1, \(\pi_{n,m}(\mathbf{C}_m)\) is a \(p\) -ring; we have \(\pi_{n,m}(\mathbf{C}_m) + \mathfrak{m}_n = \mathbf{A}_n\) , hence \(\pi_{n,m}(\mathbf{C}_m)\) is equal to the Cohen subring \(\mathbf{C}_n\) of \(\mathbf{A}_n\) . By prop. 3 of no. 1, the subring \(\varprojlim \mathbf{C}_n\) of \(\varprojlim \mathbf{A}_n\) is a \(p\) -ring. Set \(\mathbf{C} = \bigcap_{n \in \mathbb{N}} \pi_n^{-1}(\mathbf{C}_n)\) . Since C is the inverse image of \(\varprojlim \mathbf{C}_n\) under the isomorphism \(a \mapsto (\pi_n(a))_{n \in \mathbb{N}}\) from A onto \(\varprojlim \mathbf{A}_n\) , it is a closed subring of A, and a \(p\) -ring. We have \(\pi_n(\mathbf{C}) = \mathbf{C}_n\) for every \(n \in \mathbb{N}\) (no. 1, prop. 3), and in particular \(\pi_0(\mathbf{C}) = \mathbf{A}_0\) , that is, \(\pi(\mathbf{C}) = k\) . Hence C is a Cohen subring of A.

For every integer \(n \geqslant 0\) , set \(J_{n} = C \cap m_{A}^{n}\) . Since the local ring A is separated, we have \(\bigcap_{n \in \mathbb{N}} J_{n} = \{0\}\) , and in view of the structure of ideals of a \(p\) -ring (no. 1, prop. 1), every ideal of C of the form \(p^{k}C\) contains one of the \(J_{n}\) . Conversely, \(J_{n}\) contains \(p^{n}C\) . Hence, the \(pC\) -adic topology of C is induced by the \(m_{A}\) -adic topology of A. This proves \(b\) ).

Let A' be a closed subring of A, containing S and such that \(A' + m_{A} = A\). Since A' is closed, we have A' = \(\bigcap_{n\in\mathbb{N}}\pi_n^{-1}(\pi_n(A'))\) . We have \(\pi_n(A') \supset \pi_n(S)\) and \(\pi_n(A') + m_n = A_n\) , whence \(\pi_n(A') \supset C_n\) according to what was seen in A). Finally, we have \(\pi_n^{-1}(\pi_n(A')) \supset \pi_n^{-1}(C_n)\) , whence A' ⊃ C. This proves c). From this we deduce the uniqueness of a Cohen subring as in A).

Remark. --- Suppose that \(p1_{A}\) is not nilpotent (this occurs in particular when A is an integral domain whose field of fractions has characteristic 0). Then C is a discrete valuation ring whose field of fractions has characteristic 0.

